\documentclass{ximera}
%verbeterd 30/10/26
\title{Theorie: Gemengde schakelingen}
\author{Daan Lipkens}

\begin{document}
\begin{abstract}
    Inleiding tot het berekenen van stroom en spanning in schakelingen met meerdere weerstanden en één spanningsbron.
\end{abstract}
\maketitle
\label{hoofdstuk:schakelingen}
% ----------------------------------------------------------------------
%          inleiding hoofdstuk
% ----------------------------------------------------------------------
\section*{Inleiding}
Bijna elk elektrisch toestel bevat meerdere onderdelen die met elkaar verbonden zijn. Zo'n geheel van een spanningsbron, verbindingsdraden en andere componenten noemen we een \emph{elektrische schakeling}. In de praktijk zijn de componenten zelden gewoon achter elkaar gezet: er zijn vertakkingen waar de stroom zich verdeelt en weer samenkomt.

In dit hoofdstuk leer je hoe je de stroomsterkte en de spanning bepaalt in schakelingen met meerdere weerstanden, aangesloten op één spanningsbron. We starten met een korte herhaling van de grootheden en de wet van Ohm, zie~\ref{sec:herhaling-ohm}. Daarna bekijken we de twee basisschakelingen: de serieschakeling, zie~\ref{sec:serieschakeling}, en de parallelschakeling, zie~\ref{sec:parallelschakeling}.

We sluiten af met de gemengde schakeling, zie~\ref{sec:gemengde-schakeling}. Daar combineer je alles: je vervangt groepen weerstanden stap voor stap door één vervangingsweerstand en werkt daarna terug om alle stromen en spanningen te vinden.

\begin{definition}[title={Elektrische schakeling}]\nl
    Een \textbf{elektrische schakeling} is een gesloten geleidende verbinding van een spanningsbron en andere componenten, waarin een elektrische stroom kan vloeien.
\end{definition}

\end{document}